\documentclass[10pt,a4paper]{amsart}
\usepackage[margin=19mm]{geometry}
\usepackage{amsmath,amssymb,mathtools}
\usepackage[scr=rsfs,cal=euler]{mathalfa}
\usepackage{xfrac}
\usepackage[unicode,hidelinks,bookmarks=false]{hyperref}
\newcommand{\ie}{i.e.\ }
\newcommand{\cf}{cf.\ }
\newcommand{\resp}{respectively\ }
\newcommand{\Subsection}{\subsection}
\providecommand{\nobreakdash}{\nobreak\hbox{-}}
\setlength{\emergencystretch}{2em}
\multlinegap0pt
\usepackage{amssymb}
\usepackage[shortlabels]{enumitem}
\usepackage{xcolor}
\usepackage{xfrac}
\usepackage{cancel}
\newtheorem{corollary}{Corollary}
\newcommand{\N}{\mathbb{N}}
\newcommand{\Sum}{\displaystyle \sum}
\newcommand{\Z}{\mathbb{Z}}
\newcommand{\noteta}{h}
\newcommand{\pref}[1]{(\ref{#1})}
\title{Lifting Milnor Invariants for 3-Component Links}
\author{Christopher W. Davis, JungHwan Park}
\date{Source: arXiv:2605.23086v2 (CC BY 4.0)}
\begin{document}
\maketitle

\noindent\emph{Source and licence.} Lifting Milnor Invariants for 3-Component Links. Version arXiv:2605.23086v2 (CC BY 4.0), distributed by arXiv under the Creative Commons Attribution 4.0 International licence, http://creativecommons.org/licenses/by/4.0/, as recorded in the arXiv licence metadata for this exact version and shown on the abstract page https://arxiv.org/abs/2605.23086v2. Full licence notice: \texttt{licenses/arxiv-2605.23086-CC-BY-4.0.txt}.\par\smallskip

\noindent\textit{Editorial scope.} Counted unit: the article's corollary cor:lift (source0.tex lines 1587--1605). The article's complete original proof of it is delivered with it (source0.tex lines 1608--1639, one \texttt{proof} environment of 32 lines). No mathematics is written, repaired or re-derived: the statement and the proof are the article's own lines, in the article's own notation, and no retained line is substituted or re-flowed.  No further source-local context is needed: every cross-reference of every retained line resolves inside the two delivered spans.  Nothing was omitted from any retained span, so the package carries no omission record.  No retained line contains a citation, so the wrapper carries no bibliography and no bibliography entry.  No retained line contains a figure environment, an image-inclusion command or drawing code, so no image is delivered and the package declares no asset dependency.  Editorial formatting only, in addition: an AMS \texttt{amsart} preamble that restates the article's own declarations and macros used by the retained lines, namely the environments \texttt{corollary} on separate counters, together with the standard packages those lines need.  The retained environments print in this document as Corollary 1 in order of appearance, whereas the article prints the same items with the numbering of its own full document.  The complete native source of the article as distributed in the arXiv e-print for this exact version is bundled unchanged as \texttt{sources/201181/source0.tex}.
\begin{corollary}\label{cor:lift}
Let $L=(L_1, L_2, L_3)$ and $L'=(L_1',L_2', L_3')$ be weakly cobordant links, and let $k\in \N$. Assume that $\gamma^j(L)=0$ for all $j<k$ and $\gamma^k(L)\neq 0$. Then:
\begin{enumerate}
\item \label{item:cor-gamma-k}
$\gamma^k(L')=\gamma^k(L)$.
\item \label{item:cor-gamma-k+1}
$\gamma^{k+1}(L')-\gamma^{k+1}(L)=p \gamma^k(L)$ for some $p\in \Z$.
\item \label{item:cor-gamma-ell}
If $\gamma^{k+1}(L')-\gamma^{k+1}(L)=p \gamma^k(L)$ with $p\ge 0$, then 
\begin{eqnarray*}\gamma^{k+2}(L')&=&\gamma^{k+2}(L)+p\gamma^{k+1}(L)+{p\choose 2}\gamma^k(L),\\
\gamma^{k+3}(L')&=&\gamma^{k+3}(L)+p\gamma^{k+2}(L)+{p\choose 2}\gamma^{k+1}(L)+{p\choose 3}\gamma^{k}(L),
\end{eqnarray*}
and in general, if we follow the convention that $\gamma^i(L)=0$ when $i<0$ then for all $\ell>k+1$,
\begin{eqnarray*}
\gamma^\ell(L') &=& \gamma^\ell(L)+p\gamma^{\ell-1}(L)+{p\choose 2}\gamma^{\ell-2}(L)+\dots + p\gamma^{\ell-p+1}(L)+\gamma^{\ell-p}(L) \\
&=& \sum_{q=0}^p {p\choose q}\,\gamma^{\ell-q}(L).
\end{eqnarray*}
\end{enumerate}
\end{corollary}

\begin{proof}
Since $L$ and $L'$ are weakly cobordant, it follows (after the change of variables $x=t-1$) that
$$
\noteta(L')=(1+x)^n\noteta(L)
$$
for some $n\in \Z$. By swapping $L$ and $L'$ as appropriate, we may assume that $n\ge 0$. By the binomial theorem and recalling that $\gamma^j(L)=0$ for all $j<k$, we have
$$\noteta(L')=\Sum_{q=0}^n {n\choose q}x^q\cdot \Sum_{j=k}^\infty \gamma^j(L)\,x^j.$$
Expanding the product, reindexing by $j+q=m$, and using that $\gamma^{m-q}(L)=0$ when $m-q<0$, 
\begin{equation}\label{eqn: expand prod}
\noteta(L')
=
\Sum_{m=k}^\infty
\left(\Sum_{q=0}^{n}{n\choose q}\gamma^{m-q}(L)\right)x^m.
\end{equation}
We can compute the $x_k$-coefficient to get
$$
\gamma^k(L')=\Sum_{q=0}^{n}{n\choose q}\gamma^{k-q}(L)=\gamma^k(L),
$$
where the second equality follows since $\gamma^j(L)=0$ for all $j<k$.  This proves \pref{item:cor-gamma-k}.

Comparing the coefficient of $x^{k+1}$ in equation \pref{eqn: expand prod} yields
$$
\gamma^{k+1}(L')={n\choose 0}\gamma^{k+1}(L)+{n\choose 1}\gamma^{k}(L)
$$
so $\gamma^{k+1}(L')-\gamma^{k+1}(L)=n \gamma^k(L)$. This proves claim \pref{item:cor-gamma-k+1} (with $p=n$).

Finally, assume $\gamma^{k+1}(L')-\gamma^{k+1}(L)=p \gamma^k(L)$ with $p\ge 0$. Since $\gamma^k(L)\neq 0$, the previous paragraph implies that $p=n$. Comparing the coefficient of $x^\ell$ in equation~\eqref{eqn: expand prod} for any $\ell>k+1$ gives
$$
\gamma^\ell(L')=\Sum_{q=0}^{p}{p\choose q}\gamma^{\ell-q}(L),
$$
proving \pref{item:cor-gamma-ell}.
\end{proof}

\end{document}
